Large Language Models (LLMs) power modern interactive services including chatbots, virtual assistants, and programming support tools, serving millions of customer queries daily. These AI-powered services face unique operational challenges: managing stochastic customer demand, ensuring quality of service (QoS) under strict resource constraints, and balancing service capacity with responsiveness. The Key-Value (KV) cache grows dynamically during service provision, and memory overflow triggers eviction cascades that can cause system-wide failures. Even when resource capacity exceeds theoretical requirements, conventional scheduling fails because it does not account for this dynamic resource consumption---a system that should be stable can become unstable under poor scheduling.

We formulate LLM inference as a multi-stage service scheduling problem under memory constraints, where customer requests (prompts) arrive stochastically and service resource requirements (KV cache) grow dynamically during processing. Using fluid dynamics from queueing theory, we develop the Waiting for Accumulated Inference Threshold (WAIT) algorithm that prevents eviction cascades by keeping the system near load balance, achieving near-optimal service capacity while maintaining bounded customer waiting time and response latency.

For practical service settings where output lengths are unknown at arrival, we introduce Nested WAIT. Rather than predicting service times, Nested WAIT classifies customer requests on-the-fly: short requests complete early and exit, while longer requests naturally advance to later service stages. A safety buffer provides high-probability protection against resource overflow with only logarithmic overhead.

Theoretical analysis establishes near-optimal performance in the asymptotic regime. Experiments with real-world service data (LMSYS-Chat-1M) on NVIDIA A100 GPUs demonstrate 12-25\% higher service throughput and 30\% lower customer latency compared to current serving systems (vLLM, Sarathi). This work bridges service operations theory with AI systems, offering scheduling strategies directly applicable to emerging LLM-powered services.
